% Generated draft: requires mathematical and source review.
\documentclass[11pt]{article}
\usepackage[margin=25mm]{geometry}
\usepackage{fontspec}
\setmainfont{Latin Modern Roman}
\usepackage{amsmath,amssymb,mathtools,unicode-math}
\setmathfont{Latin Modern Math}
\usepackage{enumitem,needspace,xurl,hyperref}
\setlength{\parindent}{0pt}
\setlength{\parskip}{5pt}
\setlength{\emergencystretch}{3em}
\newcommand{\N}{\mathbb{N}}
\newcommand{\Z}{\mathbb{Z}}
\newcommand{\Q}{\mathbb{Q}}
\newcommand{\R}{\mathbb{R}}
\newcommand{\C}{\mathbb{C}}
\newcommand{\F}{\mathbb{F}}
\newcommand{\A}{\mathbb{A}}
\newcommand{\PP}{\mathbb{P}}
\newcommand{\E}{\mathbb{E}}
\newcommand{\eps}{\varepsilon}
\newcommand{\dd}{\,\mathrm{d}}
\begin{document}
\section*{UnsolvedMath ID: 30002676}
\textbf{Source problem code:} catalogue entry
\subsection{Definitions and mathematical statement}
Let \(N\) denote a finite positive integer counting species, and write \(c=(c_1,\ldots,c_N)\) for their state variables. Here a reactive Cahn--Hilliard system means a multicomponent evolution model combining Cahn--Hilliard diffusion with chemical reactions. A generalized gradient-flow formulation describes that evolution through an energy functional and a dissipation structure. The supplied record does not specify these objects or the equations they generate.

\textbf{Requested extension.} Extend the generalized gradient-flow formulation, well-posedness theory, and energy-dissipation theory of the source's two-component reactive Cahn--Hilliard model to reactive Cahn--Hilliard systems with an arbitrary finite number \(N\) of species. The extension should identify the applicable class of systems and hypotheses, formulate their generalized gradient-flow structure, establish well-posedness in specified solution spaces, and establish the corresponding energy-dissipation law.

\textbf{Unresolved qualifications.} This is a research objective rather than a fully specified theorem. The supplied record gives neither the reference two-component equations nor the admissible multicomponent equations. It does not specify a spatial domain, dimension, time interval, initial or boundary conditions, species constraints, reaction network, reaction rates, free energy, mobility, dissipation structure, regularity assumptions, or solution concept. Consequently, the precise quantifiers over models and data, the intended scope of well-posedness, and whether energy dissipation is expressed by an identity or an inequality remain unresolved. No particular choices of these ingredients are imposed here.
\subsection{Short English statement}
Extend the source's two-species reactive Cahn--Hilliard gradient-flow theory to arbitrarily many species, including well-posedness and energy dissipation. The supplied record leaves the underlying equations and hypotheses unspecified, so a precise theorem requires further source clarification and mathematical review.
\subsection{Sources}
\begin{itemize}
\item Supplied local UnsolvedMath registry record, permanent ID \(30002676\), source problem code \texttt{catalogue entry}, title: Gradient Flows for \(N\)-Species Reactive Cahn--Hilliard Systems. The record assigns proposed year 2014 and supplies no named proposer.
\item Supplied source citation: \emph{Reactive Flows in Deformable, Complex Media} (2015), Oberwolfach Reports, pp. 2409--2468, DOI 10.4171/owr/2014/43. \url{https://doi.org/10.4171/owr/2014/43}. The background describes Oberwolfach Reports as documenting mathematical workshops and their open-problem discussions. The citation's 2015 date and the record's 2014 report identification are retained without reconciliation.
\item Supplied literature reference: Clavijo et al., reactive \(n\)-species Cahn--Hilliard model. \url{https://doi.org/10.1016/j.cam.2018.10.007}.
\item Supplied literature reference: Extended Larch\'{e}--Cahn reactive multicomponent model. \url{https://doi.org/10.1007/s00161-021-01045-9}.
\item Supplied literature reference: Oberwolfach Report 43/2014. \url{https://doi.org/10.4171/owr/2014/43}.
\item Literature assessment dated exactly \texttt{2026-08-21}: the supplied curation labels the problem partially solved and reports that general thermodynamically consistent reactive \(n\)-species models and numerical schemes exist, while a general rigorous theory matching the requested extension remains incomplete. This is supplied dated curation derived from the OWR material, not an independently verified literature assessment; time-sensitive status and cited claims require independent verification.
\item The registry labels the statement \texttt{corrected\_verified} and reports a manual semantic, typographic, and LaTeX review describing its cleanup as meaning-preserving, self-contained, and free of unsupported source-dependent additions. These are supplied registry assertions, not independent verification of the mathematical formulation.
\end{itemize}
\end{document}
